\documentclass[10pt,twocolumn]{article}
\usepackage[margin=0.85in]{geometry}
\usepackage{graphicx,booktabs,amsmath,url,microtype}
\usepackage[hidelinks]{hyperref}
\usepackage{natbib}

\title{\bf Temporal Aggregation Erases the Advantage of\\
Surgical-Domain Pretraining}
\author{Andrew Haggstrom}
\date{}

\begin{document}
\maketitle

\begin{abstract}
Frozen vision foundation models are increasingly selected as perception
backbones for surgical video systems, and surgical-domain pretraining is
reported to confer large advantages over general-purpose encoders. We ask how
much of that advantage survives the temporal aggregation a deployed system
performs. Benchmarking four encoders spanning the pretraining
$2\times2$---general/medical $\times$ self-supervised/language-supervised---on
Cholec80 phase recognition with paired video-level statistics over all 80
videos, we find that EndoViT leads DINOv2 by 6.0 accuracy points under
frame-level linear probing ($p<0.001$), but that 85\% of this advantage
disappears under causal temporal smoothing, leaving a residual gap
statistically indistinguishable from zero ($+0.009$, $p=0.29$). The effect is
specific to the domain-pretrained encoder: every EndoViT gap shrinks by
55--85\%, whereas the DINOv2--CLIP gap is unchanged ($+0.033 \to +0.032$, both
$p<0.001$), ruling out generic compression from smoothing. Our aggregator is a
majority filter, substantially weaker than the multi-stage temporal
convolutional networks used in practice, so this constitutes a lower bound on
what a learned temporal head would recover. We additionally find that medical
pretraining delivered by language supervision transfers essentially nothing
(BiomedCLIP $\approx$ CLIP at every window, a well-powered null), that all four
encoders over-segment surgical phase by 80--124$\times$ at frame accuracies of
0.66--0.72, and that raw inter-frame drift is confounded by feature variance
and must be normalized. We further report a result we retracted: an apparent
rank inversion observed on a 32-video test split did not replicate at $n=80$,
which bears on how the sub-two-point encoder margins common in this literature
should be read. Code, cached embeddings, and evaluation protocol are released.
\end{abstract}

\section{Introduction}

Frozen vision foundation models are now a default perception backbone for
surgical video systems. The practical decision a group faces is which backbone
to adopt, and increasingly whether to pay for surgical-domain pretraining
\citep{batic2024endovit,jaspers2025surgenetxl,zen2026} over a strong
general-purpose encoder such as DINOv2 \citep{oquab2024dinov2}.

That decision is made empirically, and the answer depends on the evaluation
regime. Two coexist. The first is \emph{frame-level probing}: linear probes and
$k$-nearest-neighbour classifiers fit to frozen features from independently
sampled frames, reported throughout the representation-learning literature as
proxies for representation quality \citep{ramesh2023dissecting,surglavi2026}.
The second is \emph{frozen-backbone with temporal aggregation}: the two-stage
protocol of TeCNO \citep{czempiel2020tecno}, in which cached frame-level
features are passed to a multi-stage temporal convolutional network
\citep{farha2019mstcn}. This is the protocol used for the frozen-backbone
setting in recent surgical foundation model evaluations \citep{zen2026,msics2026}.

These are not the same measurement. The first is what most papers report when
comparing encoders; the second is what a deployed system instantiates. The
relationship between them---how much of a frame-level advantage survives
aggregation, and whether that survival depends on how the encoder was
pretrained---is, to our knowledge, unmeasured.

We measure it. Our contributions:

\begin{enumerate}\itemsep2pt
\item \textbf{The domain-pretraining advantage is largely
temporal-recoverable.} EndoViT's 6.0-point frame-level lead over DINOv2
($p<0.001$, 80 paired videos) shrinks by 85\% under causal smoothing to a
non-significant $+0.009$.
\item \textbf{The effect is specific to the domain-pretrained encoder.} All
EndoViT gaps shrink 55--85\%; the DINOv2--CLIP gap does not move. This
localizes the effect to pretraining type rather than to smoothing itself.
\item \textbf{Supervision matters more than domain.} BiomedCLIP is
indistinguishable from CLIP at every window; both self-supervised encoders beat
both language-supervised ones.
\item \textbf{Magnitude.} All four encoders over-segment surgical phase by
80--124$\times$ at frame accuracies that read as respectable.
\item \textbf{A lower bound.} A majority filter is far weaker than MS-TCN; that
it already erases 85\% bounds from below what a learned head recovers.
\item \textbf{A retracted result} and its implication for how small encoder
margins in this literature should be interpreted (\S\ref{sec:retraction}).
\end{enumerate}

\section{Related work}

\paragraph{Frozen-feature evaluation and its fragility.}
Concerns about probe-based model comparison predate this application:
\citet{resnick2019probing} showed that models ranking poorly under linear
probing can rank better under alternative readouts. Our observation that linear
probing and $k$-NN select different encoders (\S\ref{sec:frame}) is a surgical
instance of this known phenomenon and serves as motivation rather than as a
novel finding. The reproducibility dimension has been documented sharply
outside medicine: an audit of geospatial foundation models \citep{gfm2026sota}
finds that papers reporting the same checkpoint on the same benchmark under the
same nominal linear-probe protocol disagree by up to 12.7 points at the
90th percentile, largely because probe recipes go undisclosed.

\paragraph{Temporal models for surgical phase recognition.}
The line from EndoNet \citep{twinanda2017endonet} through TeCNO
\citep{czempiel2020tecno}, Trans-SVNet \citep{gao2021transsvnet}, LoViT
\citep{liu2025lovit} and SKiT \citep{liu2023skit} establishes that frame-level
features require temporal modelling to reach usable accuracy. These
architectures exist \emph{because} frozen frame-level features are temporally
unreliable. We quantify what their aggregation recovers and show that this
recovery substantially substitutes for domain-specific pretraining. Supporting
evidence already appears without being named as such: \citet{surglavi2026}
report that a linear probe on frozen features comes within 2.34\% F1 of a
specialised architecture, and that jointly optimised backbone-plus-aggregation
models retain only a 7.4\% average advantage despite far greater complexity.

\paragraph{Surgical and medical foundation models.}
EndoViT \citep{batic2024endovit} applies masked autoencoding \citep{he2022mae}
to Endo700k; SurgeNetXL \citep{jaspers2025surgenetxl} scales surgical SSL;
BiomedCLIP \citep{zhang2023biomedclip} provides medical-domain pretraining via
language supervision on biomedical figure--caption pairs. Reported gains are
protocol-dependent in a manner consistent with our finding:
\citet{jaspers2025surgenetxl} report procedure-specific pretraining improving
Dice by 10.0--29.6\% over ImageNet initialisation with \emph{frozen} weights but
only 2.8--12.8\% with trainable weights---the frozen setting inflates the
apparent benefit roughly two- to threefold.

\paragraph{Closest prior work.}
A controlled study of phase segmentation in small-incision cataract surgery
\citep{msics2026} uses a shared MS-TCN++ head over frozen cached features across
several backbones, and cautions that deployment decisions should not rely on
peak accuracy alone, noting that domain pretraining yields mixed,
backbone-dependent transfer. That study compares backbones \emph{under} a fixed
temporal head; we compare the same encoders \emph{across} the frame-level and
aggregated regimes and measure the change in each pairwise gap, with paired
video-level intervals that distinguish a shrinking-but-real advantage from one
that has become indistinguishable from zero.

\section{Method}

\subsection{Encoders}
We evaluate four encoders spanning the pretraining $2\times2$
(Table~\ref{tab:encoders}): DINOv2 ViT-S/14 \citep{oquab2024dinov2} (general,
self-supervised), CLIP ViT-B/16 \citep{radford2021clip} (general,
language-supervised), EndoViT ViT-B/16 \citep{batic2024endovit} (surgical,
self-supervised MAE), and BiomedCLIP ViT-B/16 \citep{zhang2023biomedclip}
(medical, language-supervised).

All encoders are frozen and used in evaluation mode. Each applies its own
normalisation statistics---ImageNet, CLIP, endoscopic, and PMC respectively---as
using incorrect statistics silently degrades features. Frames are resized to
$224\times224$ with bicubic interpolation and antialiasing. We take the
pre-projection pooled representation for both CLIP variants so that all four
encoders are probed on comparable features. Implementations use \texttt{timm}
\citep{wightman2019timm} and \texttt{open\_clip} \citep{ilharco2021openclip}.

Embeddings are extracted once and cached as memory-mapped arrays partitioned by
a hash of the encoder configuration, with SHA-256 manifests verified on every
read; all downstream experiments are pure functions over this cache.

\subsection{Evaluation protocols}
\textbf{P1, linear probe}: class-balanced multinomial logistic regression on
standardised frozen features. \textbf{P2, $k$-NN}: $k=20$, cosine metric, no
learned head. \textbf{P3, temporally smoothed accuracy}: a majority filter of
window $w$ applied to the P1 prediction stream. We report both \emph{centered}
(offline) and \emph{causal/trailing} (real-time; the only deployable variant)
filters; $w=1$ recovers P1 exactly.

\subsection{Temporal reliability metrics}
All metrics are computed within a single contiguous sequence and never across
video boundaries. \emph{Normalized drift} is the mean inter-frame cosine
distance divided by the sequence's spread about its centroid; the normalisation
is necessary rather than cosmetic (\S\ref{sec:variance}). \emph{Boundary
jitter} is the predicted phase-transition rate as a ratio to the ground-truth
rate. \emph{Conditional stability} multiplies stability by neighbour-label
consistency, penalising the degenerate constant-output solution.
\emph{Velocity--flow coupling} is the Spearman correlation between embedding
velocity and dense Farneback optical-flow magnitude, separating drift that
tracks scene motion from drift that does not. Twelve synthetic controls with
known answers gate the suite.

\subsection{Data and statistics}
We use Cholec80 \citep{twinanda2017endonet}: 80 laparoscopic cholecystectomy
videos with per-frame phase annotations over seven phases and binary presence
labels for seven instruments. Following the standard protocol we sample at
1\,fps, yielding 184{,}498 frames. Phase labels are provided at 25\,fps and
instrument labels at 1\,fps; we align them by taking the phase label at each
sampled native frame index and assert index agreement rather than assuming it.
Pilot experiments used CholecSeg8k \citep{hong2020cholecseg8k}.

Frames within a surgery are near-duplicates, so \textbf{all splitting and
resampling is at the video level}. Our primary analysis is 5-fold
cross-validation over all 80 videos with a shared fold assignment across
encoders, so that every comparison is paired; each video is scored exactly once
by a probe trained on the other four folds. Pairwise differences are assessed
by a paired video-level bootstrap \citep{efron1994bootstrap} with 10{,}000
resamples of the per-video difference. Pairing cancels the large between-video
variation in duration and difficulty; overlapping marginal intervals do not
imply a non-significant paired difference. For comparability with published
work we also report the fixed 40/8/32 split. Phase prevalence varies by up to
$31\times$, so macro-F1 is reported alongside accuracy.

Our panel comprises 6 pairs $\times$ 5 windows $=$ 30 tests, of which 24 are
significant at $p<0.05$ against $\approx$1.5 expected by chance; the null
results below are therefore informative rather than merely underpowered.

\begin{table}[t]\centering\small
\caption{Encoders. VRAM is measured peak allocation during extraction.
Throughput is I/O-bound in our pipeline and is not model throughput.}
\label{tab:encoders}
\begin{tabular}{@{}llrr@{}}
\toprule
Encoder & Pretraining & Dim & VRAM \\
\midrule
DINOv2 ViT-S/14   & LVD-142M, SSL      & 384 & 489\,MB \\
CLIP ViT-B/16     & WIT-400M, language & 768 & 1011\,MB \\
EndoViT ViT-B/16  & Endo700k, MAE      & 768 & 856\,MB \\
BiomedCLIP ViT-B/16 & PMC-15M, language & 768 & 856\,MB \\
\bottomrule
\end{tabular}
\end{table}

\section{Results}

\subsection{Frame-level protocols disagree}
\label{sec:frame}

\begin{table}[t]\centering\small
\caption{Frame-level protocols, fixed 40/8/32 split. P1 selects EndoViT;
P2 selects BiomedCLIP.}
\label{tab:frame}
\begin{tabular}{@{}lrrrrr@{}}
\toprule
& \multicolumn{2}{c}{P1 linear} & \multicolumn{2}{c}{P2 $k$-NN} & \\
Encoder & acc & F1 & acc & F1 & gap \\
\midrule
EndoViT    & \bf.704 & \bf.615 & .522 & .377 & .238 \\
DINOv2     & .677 & .597 & .611 & .463 & .134 \\
BiomedCLIP & .658 & .578 & \bf.617 & \bf.504 & .075 \\
CLIP       & .648 & .568 & .572 & .391 & .177 \\
\bottomrule
\end{tabular}
\end{table}

Table~\ref{tab:frame} shows two standard protocols selecting different winners.
EndoViT has the best linear separability and the worst geometric clustering: its
phases are separable by a hyperplane but not by proximity, consistent with it
having the smallest feature spread of the four (embedding standard deviation
0.425 versus 2.46 for DINOv2).

The per-phase pattern is identical across all four encoders: the long
dissection phases score 0.69--0.77 F1 while transitional and rare phases
collapse to 0.46--0.59. This failure is structural rather than model-specific,
and concerns exactly the states a single frame cannot disambiguate.

\subsection{Temporal reliability at $w=1$}
Every encoder over-segments dramatically. Predicted phase-transition rates
exceed ground truth by $79.7\times$ (EndoViT), $110.5\times$ (DINOv2),
$122.4\times$ (BiomedCLIP) and $124.3\times$ (CLIP), at frame accuracies of
0.66--0.72. The best encoder---a surgical foundation model pretrained on this
data distribution---still announces eighty times too many phase transitions.
This is a within-encoder measurement against ground truth, two orders of
magnitude in size, and does not depend on any cross-encoder comparison.

\subsection{The domain advantage converges away}
\label{sec:converge}

\begin{table}[t]\centering\small
\caption{Mean accuracy, 5-fold CV over all 80 videos, causal smoothing.
Ordering is preserved at every window; the \emph{gaps} are what change.}
\label{tab:acc}
\begin{tabular}{@{}lrrrrr@{}}
\toprule
Encoder & $w{=}1$ & 15 & 31 & 61 & 121 \\
\midrule
EndoViT    & \bf.756 & \bf.805 & \bf.810 & \bf.795 & \bf.751 \\
DINOv2     & .697 & .779 & .790 & .784 & .742 \\
BiomedCLIP & .676 & .760 & .770 & .762 & .723 \\
CLIP       & .664 & .753 & .764 & .752 & .710 \\
\bottomrule
\end{tabular}
\end{table}

\begin{table}[t]\centering\small
\caption{Paired accuracy differences, causal smoothing, $n=80$ videos.
Shrinkage is $1-\Delta_{121}/\Delta_{1}$. All EndoViT gaps shrink; no
general-encoder gap does.
$^{***}p{<}.001$, $^{**}p{<}.01$, $^{*}p{<}.05$, n.s.\ otherwise.}
\label{tab:paired}
\begin{tabular}{@{}lrrrr@{}}
\toprule
Pair & $w{=}1$ & $w{=}31$ & $w{=}121$ & shrink \\
\midrule
EndoViT$-$DINOv2     & $+.060^{***}$ & $+.019^{*}$   & $+.009$       & \bf85\% \\
EndoViT$-$BiomedCLIP & $+.080^{***}$ & $+.040^{***}$ & $+.029^{**}$  & 64\% \\
EndoViT$-$CLIP       & $+.092^{***}$ & $+.046^{***}$ & $+.041^{***}$ & 55\% \\
DINOv2$-$CLIP        & $+.033^{***}$ & $+.026^{***}$ & $+.032^{***}$ & \bf2\% \\
DINOv2$-$BiomedCLIP  & $+.021^{**}$  & $+.020^{*}$   & $+.020^{*}$   & 5\% \\
BiomedCLIP$-$CLIP    & $+.012^{*}$   & $+.006$       & $+.013$       & --- \\
\bottomrule
\end{tabular}
\end{table}

Table~\ref{tab:paired} is the paper's central result. EndoViT's advantage over
DINOv2 crosses from significant to non-significant between $w=31$ ($p=0.017$)
and $w=61$ ($p=0.19$); at $w=121$ EndoViT wins on 44 of 80 videos, a coin flip.
Every EndoViT gap shrinks by 55--85\%. No gap between general-purpose encoders
shrinks at all: DINOv2$-$CLIP is $+0.033$ at $w=1$ and $+0.032$ at $w=121$, both
$p<0.001$.

This contrast is what makes the effect attributable to pretraining type. If
smoothing merely compressed all differences toward a ceiling, the
general-encoder gaps would shrink too. They do not. The interpretation we
favour is that surgical pretraining principally buys single-frame
discriminability, and that temporal aggregation recovers much of the same
information from context---so the two routes to the same signal overlap.

Accuracy peaks near $w\approx31$ under causal smoothing and declines thereafter,
because a trailing window necessarily lags every transition; macro-F1 peaks
near $w\approx15$--31 and falls below the unsmoothed baseline by $w=121$ as
short rare phases are erased. We therefore take $w\approx31$ as the honest
operating point. Under centered smoothing accuracy rises monotonically---an
artifact of hindsight that the causal filter exposes.

Across the sweep, jitter falls 38--54$\times$ while accuracy \emph{rises}: the
phase information was present in the frozen features throughout, and P1 gave no
indication of how much temporal repair each encoder would require.

\subsection{Supervision beats domain}
BiomedCLIP and CLIP are statistically indistinguishable at every window
$w\geq15$ ($p=0.17$--$0.47$, differences $\leq0.013$, splits 33--42 of 80).
Medical-domain pretraining delivered through language supervision transfers
essentially nothing to surgical video, whereas medical-domain self-supervised
pretraining yields a large frame-level lead. Both self-supervised encoders beat
both language-supervised encoders at every window. A plausible mechanism is
that PMC-15M comprises static biomedical figures---histology, radiology,
diagrams---whose visual statistics differ sharply from endoscopic video;
nominal domain proximity does not imply transfer. Given that 24 of 30
comparisons in the same panel reach significance, this is a well-powered null.

\subsection{Is drift signal or noise?}
Spearman correlations between embedding velocity and optical-flow magnitude
(10 test videos, all $p<0.05$) are 0.499 (EndoViT), 0.471 (DINOv2) and 0.388
(CLIP). Drift partly tracks scene motion, but roughly half the variance is
unexplained by pixel motion, indicating a substantial spurious component. The
ordering matches the stability ranking.

\section{Analysis}

\subsection{Raw drift is variance-gameable}
\label{sec:variance}
An early run indicated CLIP as \emph{more} temporally stable than DINOv2 (raw
inter-frame drift 0.042 versus 0.124). Diagnosis: CLIP's feature spread was
$3.5\times$ smaller and its embedding standard deviation $2.7\times$ smaller.
Its drift was low because its features were low-variance, not because they were
stable. The internal contradiction that flagged this was that CLIP
simultaneously jittered \emph{more} ($124\times$ versus $110\times$).

Normalising drift by sequence spread removes the artifact. Notably EndoViT has
the \emph{lowest} raw variance of the four (0.425) yet is genuinely the most
stable once normalised (drift ratio 0.512 versus 0.665 and 0.800)---low
variance realised well rather than badly. Raw drift and raw variance each
mislead in isolation; only the normalised metric distinguishes a bland encoder
from a well-organised one.

\subsection{A retracted result}
\label{sec:retraction}
On the fixed 40/8/32 split (32 test videos, 40 training), causal smoothing
appeared to \emph{invert} the DINOv2/EndoViT ordering: EndoViT $+0.030$ at
$w=1$, DINOv2 $+0.021$ at $w=121$. We initially treated this as the paper's
principal finding.

It did not replicate. Under 5-fold cross-validation (80 test videos, 64
training) the sign never flips and EndoViT leads at every window. On the
32-video split, no pairwise comparison in the panel reached significance,
including EndoViT's frame-level lead ($p=0.076$). Both the evaluation-set size
and the training-set size differ between the two analyses, so we do not
attribute the discrepancy to either alone.

The implication extends beyond our own error: sub-two-point margins between
frozen encoders on a 32-video Cholec80 test split are not resolvable, and such
rankings are routinely reported in this literature without paired error bars.
We recommend paired video-level bootstrapping and cross-validation over fixed
splits whenever a reported margin is small.

\section{Limitations}
\textbf{Domain overlap.} Endo700k includes Cholec80, so EndoViT observed these
videos during self-supervised pretraining. This inflates its frame-level lead
and therefore makes the convergence result conservative. \textbf{Single
dataset.} Cross-procedure replication, for which AutoLaparo
\citep{wang2022autolaparo} is the natural target, remains outstanding.
\textbf{Aggregator strength.} A majority filter is far weaker than MS-TCN
\citep{farha2019mstcn}; our result is a lower bound, and a learned head would
plausibly erase more of the domain advantage. \textbf{Panel size.} Four
encoders suffice for pairwise claims at $n=80$ videos but not for rank
correlations. \textbf{Efficiency.} Reported throughput is pipeline-level and
I/O-bound, not model-level.

\section{Conclusion}
Surgical-domain pretraining confers a large, unambiguous advantage on the
frame-level probe the field uses to compare frozen backbones---6.0 accuracy
points for EndoViT over DINOv2 across 80 paired videos. Under temporal
aggregation, 85\% of that advantage disappears and the remainder is
indistinguishable from zero, while gaps between general-purpose encoders are
untouched. Meanwhile every encoder tested over-segments surgical phase by
80--124$\times$ at frame accuracies that read as respectable. Practitioners
should evaluate frozen backbones under the aggregation regime they intend to
deploy, report paired error bars on any margin below a few points, and treat
frame-level gains from domain pretraining as an upper bound on deployed
benefit.

\bibliographystyle{plainnat}
\bibliography{refs}
\end{document}
